\providecommand{\Res}[2]{\qty[round-mode=figures,round-precision=3]{\fpeval{#1}}{#2}}
\providecommand{\Val}[1]{\num[round-mode=figures,round-precision=4]{\fpeval{#1}}}
\providecommand{\ResSci}[2]{\qty[round-mode=figures,round-precision=3,exponent-mode=scientific]{\fpeval{#1}}{#2}}
\subsubsection*{Question a) Intensité du courant}
\begin{enumerate}[label=\arabic*)]
    \item \textbf{Variables connues et recherchées :} $P = \qty{\v{P}}{W}$, $U = \qty{240}{V}$ ; variable recherchée : $I$
    \item \textbf{Équations pertinentes :}
    \[ P = UI \]
    \item \textbf{Isolation de la variable recherchée :}
    \[ I = \frac{P}{U} \]
    \item \textbf{Résolution numérique :}
    \[ I = \frac{\qty{\v{P}}{W}}{\qty{240}{V}} = \Res{\v{P}/240}{A} \]
    \item \textbf{Réponse :} le courant est de \textbf{\Res{\v{P}/240}{A}}.
\end{enumerate}

\subsubsection*{Question b) Résistance}
\begin{enumerate}[label=\arabic*)]
    \item \textbf{Variables connues et recherchées :} $U = \qty{240}{V}$, $I = \Res{\v{P}/240}{A}$ ; variable recherchée : $R$
    \item \textbf{Équations pertinentes :}
    \[ U = RI \]
    \item \textbf{Isolation de la variable recherchée :}
    \[ R = \frac{U}{I} \]
    \item \textbf{Résolution numérique :}
    \[ R = \frac{\qty{240}{V}}{\Val{\v{P}/240}\,\unit{A}} = \Res{240^2/\v{P}}{\ohm} \]
    \item \textbf{Réponse :} l'élément chauffant a une résistance de \textbf{\Res{240^2/\v{P}}{\ohm}}.
\end{enumerate}

\subsubsection*{Question c) Énergie consommée}
\begin{enumerate}[label=\arabic*)]
    \item \textbf{Variables connues et recherchées :} $P = \Res{\v{P}/1000}{kW}$, $\Delta t = 30 \times \qty{\v{h}}{h} = \Res{30*\v{h}}{h}$ ; variable recherchée : $E$
    \item \textbf{Équations pertinentes :}
    \[ E = P\,\Delta t \]
    \item \textbf{Isolation de la variable recherchée :} l'équation est déjà isolée.
    \item \textbf{Résolution numérique :}
    \[ E = \Val{\v{P}/1000}\,\unit{kW} \cdot \Val{30*\v{h}}\,\unit{h} = \Res{\v{P}*\v{h}*30/1000}{\kilo\watt\hour} \]
    En joules ($\qty{1}{\kilo\watt\hour} = \qty{3.6e6}{J}$) : $E = \ResSci{\v{P}*\v{h}*30*3600}{J}$.
    \item \textbf{Réponse :} $E = \mathbf{\Res{\v{P}*\v{h}*30/1000}{\kilo\watt\hour}} = \mathbf{\ResSci{\v{P}*\v{h}*30*3600}{J}}$.
\end{enumerate}

\subsubsection*{Question d) Coût}
\[ \text{Coût} = \Val{\v{P}*\v{h}*30/1000}\,\unit{\kilo\watt\hour} \times \qty{\v{tarif}}{\$\per\kilo\watt\hour} = \Res{\v{P}*\v{h}*30/1000*\v{tarif}}{\$} \]
Le chauffage de la pièce coûte \textbf{\Res{\v{P}*\v{h}*30/1000*\v{tarif}}{\$}} pour le mois.
